% 主旨：MOBA 强化学习依赖大规模训练；动作掩码在采样前施加才与策略梯度相容；
%       塑形奖励只在势函数形式下保证最优策略不变。本文落在三者交叉的短训场景。
\section{相关工作}
\label{sec:related}

\paragraph{MOBA 游戏中的强化学习。}
OpenAI Five 在 Dota 2 上以长时间大规模自对弈训练击败了职业队伍\cite{berner2019dota}。在王者荣耀上，\textcite{ye2020mastering} 针对 1v1 提出解耦的多分支动作、目标注意力和双裁剪 PPO，\textcite{ye2020full} 将其扩展到包含 \num{40} 名英雄的 5v5 对战，\textcite{wei2022hok} 公开了 1v1 环境与基线\cite{hoksource}。这些工作说明实体编码加 LSTM 的 PPO 框架足以学会对线与推塔，但其结论建立在充足的训练量之上；课程平台的基线沿用同一框架\cite{kaiwu63}，在 \qty{30}{\minute} 量级的预算下能否学到推塔行为，上述工作没有回答。

\paragraph{无效动作掩码。}
在离散动作空间中，把不合法动作的对数几率置为负无穷后再做 softmax，是处理动态合法集合的常用做法；\textcite{ye2020mastering} 也用游戏知识剪枝动作。\textcite{huang2022masking} 证明，这种掩码对应于对重新归一化后分布求出的有效策略梯度，并在实验中显示它优于对无效动作施加惩罚。由此可知，采样与梯度计算必须使用同一掩码；若在采样之后替换动作，记录的概率就不再描述实际行为。上述无效动作掩码主要处理环境层面的合法性；本文在同一机制中进一步施加针对特定英雄的战术约束，并保证掩码在 PPO 的新旧策略间一致。

\paragraph{奖励塑形。}
势函数塑形奖励 $F=\gamma\Phi(s')-\Phi(s)$ 在折扣与终止条件一致时不改变最优策略\cite{ng1999}，因而常被用来加速稀疏奖励任务。本文的差分项形式上接近势函数差，但不满足这一条件（\cref{sec:reward}），整套混合奖励不保证最优策略不变；因此本文以固定对手上的胜负而非训练回报评价策略，避免把回报上升误当作能力提升。
